% =====================================================================
%  riesgos.tex  --  Lista de Riesgos de LEIP
%  ARCHIVO UNICO: no necesita ningun otro archivo.
%  Compilar en Overleaf con pdfLaTeX (dos pasadas por el indice).
% =====================================================================
\documentclass[11pt,a4paper]{article}

% ---------------------------------------------------------------------
%  DATOS DEL PROYECTO
% ---------------------------------------------------------------------
\newcommand{\proyectoNombre}{LEIP}
\newcommand{\proyectoNombreLargo}{LatAm Economic Intelligence Platform}
\newcommand{\proyectoUniversidad}{Universidad Cat\'olica del Uruguay}
\newcommand{\proyectoFacultad}{Facultad de Ingenier\'ia y Tecnolog\'ias}
\newcommand{\proyectoCarrera}{Ingenier\'ia en Inform\'atica}
\newcommand{\docVersion}{1.1}
\newcommand{\docTitulo}{Lista de Riesgos}
\newcommand{\docFecha}{Septiembre 2026}

% ---------------------------------------------------------------------
%  PREAMBULO
% ---------------------------------------------------------------------
\usepackage[utf8]{inputenc}
\usepackage[T1]{fontenc}
\IfFileExists{lmodern.sty}{\usepackage{lmodern}\usepackage{microtype}}{}

\IfFileExists{spanish.ldf}{%
  \usepackage[spanish,es-tabla,es-noquoting]{babel}%
}{%
  \usepackage[english]{babel}%
  \renewcommand{\contentsname}{\'Indice}%
  \renewcommand{\listfigurename}{\'Indice de figuras}%
  \renewcommand{\listtablename}{\'Indice de tablas}%
  \renewcommand{\tablename}{Tabla}%
  \renewcommand{\figurename}{Figura}%
}

\usepackage[a4paper,margin=2.5cm,headheight=15pt]{geometry}
\usepackage{fancyhdr}
\usepackage{parskip}
\usepackage{booktabs}
\usepackage{tabularx}
\usepackage{longtable}
\usepackage{array}
\usepackage{enumitem}
\setlist{noitemsep,topsep=4pt}
\usepackage{graphicx}
\usepackage{xcolor}
\usepackage{colortbl}
\usepackage[hidelinks]{hyperref}
\usepackage{url}
\urlstyle{same}

\definecolor{leipAzul}{HTML}{1F4E79}
\definecolor{leipGris}{HTML}{5A5A5A}

\pagestyle{fancy}
\fancyhf{}
\fancyhead[L]{\small\color{leipGris}\proyectoNombre{} --- \docTitulo}
\fancyhead[R]{\small\color{leipGris} v\docVersion}
\fancyfoot[C]{\small\thepage}
\renewcommand{\headrulewidth}{0.4pt}

\usepackage{titlesec}
\titleformat{\section}{\Large\bfseries\color{leipAzul}}{\thesection}{0.6em}{}
\titleformat{\subsection}{\large\bfseries\color{leipAzul}}{\thesubsection}{0.6em}{}
\titleformat{\subsubsection}{\normalsize\bfseries\color{leipGris}}{\thesubsubsection}{0.6em}{}

% --- Portada -----------------------------------------------------------
\newcommand{\logoUniversidad}{img/logo-ucu.png}
\newcommand{\portada}[1]{%
  \begin{titlepage}
  \centering
  \vspace*{1.5cm}
  \IfFileExists{\logoUniversidad}{%
    \includegraphics[width=0.40\textwidth]{\logoUniversidad}\par
    \vspace{1.2cm}%
  }{%
    \fbox{\parbox[c][3cm][c]{0.5\textwidth}{\centering\small\texttt{\logoUniversidad}}}\par
    \vspace{1.2cm}%
  }
  {\Large \proyectoUniversidad\par}
  {\large \proyectoFacultad\par}
  {\large \proyectoCarrera\par}
  \vspace{3cm}
  {\Huge\bfseries\color{leipAzul} #1\par}
  \vspace{0.8cm}
  {\Large \proyectoNombre{} --- \proyectoNombreLargo\par}
  \vspace{2cm}
  {\large Versi\'on \docVersion\par}
  {\large \docFecha\par}
  \vfill
  \begin{tabular}{c}
    Sebastian Forische \\
    Armando Hernandez de la Vega \\
    Luana Acu\~na \\
    Bruno Sebastian Olivera Viera Gomez \\
  \end{tabular}
  \vspace{1.5cm}
  \end{titlepage}
}

% --- Historial de revisiones -------------------------------------------
\newenvironment{historial}{%
  \section*{Historial de revisiones}
  \addcontentsline{toc}{section}{Historial de revisiones}
  \begin{tabular}{@{}l l p{0.52\textwidth} l@{}}
  \toprule
  \textbf{Versi\'on} & \textbf{Fecha} & \textbf{Descripci\'on} & \textbf{Autor} \\
  \midrule
}{%
  \bottomrule
  \end{tabular}
  \bigskip
}

% =====================================================================
%  ENTORNO PARA DESCRIBIR UN RIESGO
% ---------------------------------------------------------------------
%  Uso:
%
%  \begin{riesgo}{Titulo del riesgo}
%    \rdescripcion{En que consiste}
%    \rcausa{Que lo origina}
%    \rimpacto{Que consecuencia tiene si ocurre}
%    \rprobabilidad{}      % segun la escala de la seccion 2
%    \rnivel{}             % probabilidad x impacto
%    \rmitigacion{Que se hace para reducirlo}
%    \rcontingencia{Que se hace si igual ocurre}
%    \rresponsable{}
%  \end{riesgo}
%
%  El codigo (R-01, R-02, ...) se numera solo.
% =====================================================================

\newcounter{rcounter}
\newcommand{\dosdig}[1]{\ifnum\value{#1}<10 0\fi\arabic{#1}}

\newcommand{\rcampo}[2]{%
  \par\noindent
  \begin{tabularx}{\textwidth}{@{}>{\bfseries}p{3.4cm} X@{}}
    #1 & #2
  \end{tabularx}\par\smallskip
}

\newcommand{\rdescripcion}[1]{\rcampo{Descripci\'on}{#1}}
\newcommand{\rcausa}[1]{\rcampo{Causa}{#1}}
\newcommand{\rimpacto}[1]{\rcampo{Impacto}{#1}}
\newcommand{\rprobabilidad}[1]{\rcampo{Probabilidad}{#1}}
\newcommand{\rnivel}[1]{\rcampo{Nivel}{#1}}
\newcommand{\rmitigacion}[1]{\rcampo{Mitigaci\'on}{#1}}
\newcommand{\rcontingencia}[1]{\rcampo{Contingencia}{#1}}
\newcommand{\rresponsable}[1]{\rcampo{Responsable}{#1}}

\newenvironment{riesgo}[1]{%
  \refstepcounter{rcounter}%
  \subsection*{\color{leipAzul}R-\dosdig{rcounter}\quad #1}
  \addcontentsline{toc}{subsection}{R-\dosdig{rcounter}\quad #1}
  \vspace{-0.4em}
  {\color{leipGris}\hrule height 0.5pt}
  \smallskip
}{%
  \bigskip
}

\begin{document}

\portada{\docTitulo}

\tableofcontents
\bigskip

\begin{historial}
1.0 & 26/06/2026 & Contenido inicial derivado del anteproyecto & Equipo \\
1.1 & 11/09/2026 & Se completan objetivo, alcance, escalas de evaluaci\'on, matriz y seguimiento. Se ampl\'ia el registro a nueve riesgos, uno por cada asunci\'on del Documento de Visi\'on & Equipo \\
\end{historial}
\clearpage

% =====================================================================
\section{Introducci\'on}

\subsection{Objetivo del documento}

Este documento registra los riesgos que pueden afectar la construcci\'on y la
operaci\'on de \proyectoNombre{}, y para cada uno deja asentados su causa, su
impacto, su nivel de exposici\'on, la acci\'on prevista para reducirlo y la
acci\'on prevista para el caso de que ocurra de todos modos.

Es un registro vivo, no un inventario cerrado: se revisa peri\'odicamente, los
riesgos cambian de nivel a medida que avanza el trabajo, y los que dejan de ser
relevantes se marcan como cerrados en lugar de eliminarse. Cada cambio queda
asentado en el historial de revisiones.

\subsection{Alcance}

El registro cubre los riesgos t\'ecnicos, de datos, legales, de dependencias
externas y de adopci\'on que afectan directamente al producto. Cada asunci\'on
enunciada en el Documento de Visi\'on tiene aqu\'i su riesgo correspondiente:
si la asunci\'on no se cumple, el riesgo asociado se materializa.

No cubre los riesgos de gesti\'on del trabajo ---disponibilidad, cronograma o
coordinaci\'on--- que se tratan en el Plan de Proyecto, ni los riesgos
comerciales del modelo de negocio, que se tratan en el Documento de Negocio.

% =====================================================================
\section{Par\'ametros de evaluaci\'on}

Cada riesgo se eval\'ua con dos variables independientes ---la probabilidad de
que ocurra y el impacto que tendr\'ia si ocurriera--- que se combinan en un
\'unico nivel de exposici\'on. Las tres escalas se definen a continuaci\'on.

\subsection{Escala de probabilidad}

La probabilidad expresa cu\'an esperable es que el evento ocurra durante el
per\'iodo de construcci\'on y operaci\'on inicial del producto.

\begin{tabularx}{\textwidth}{@{}c l X@{}}
\toprule
\textbf{Valor} & \textbf{Nivel} & \textbf{Criterio} \\
\midrule
1 & Muy baja & No hay antecedentes ni indicios de que pueda ocurrir. \\
2 & Baja     & Es concebible, pero no se espera que ocurra. \\
3 & Media    & Puede ocurrir; hay antecedentes en contextos comparables. \\
4 & Alta     & Es esperable que ocurra al menos una vez. \\
5 & Muy alta & Ocurre de forma habitual o ya se observ\'o durante el trabajo. \\
\bottomrule
\end{tabularx}

\subsection{Escala de impacto}

El impacto expresa la consecuencia sobre el producto si el evento se
materializa, medida por el esfuerzo de recuperaci\'on y por el efecto sobre la
informaci\'on que el producto entrega.

\begin{tabularx}{\textwidth}{@{}c l X@{}}
\toprule
\textbf{Valor} & \textbf{Nivel} & \textbf{Criterio} \\
\midrule
1 & Muy bajo & Se resuelve dentro del trabajo corriente, sin efecto visible. \\
2 & Bajo     & Requiere trabajo adicional acotado; el producto sigue operando. \\
3 & Medio    & Degrada una capacidad o interrumpe temporalmente una fuente. \\
4 & Alto     & Compromete la confiabilidad del dato entregado o bloquea un m\'odulo completo. \\
5 & Muy alto & Obliga a retirar informaci\'on ya publicada o a redefinir el alcance. \\
\bottomrule
\end{tabularx}

\subsection{C\'alculo del nivel de riesgo}

El nivel de riesgo resulta de multiplicar probabilidad por impacto, lo que da un
valor entre 1 y 25. Ese valor se traduce en una de cuatro categor\'ias, que
determinan c\'omo se trata el riesgo.

\begin{tabularx}{\textwidth}{@{}c l X@{}}
\toprule
\textbf{Valor} & \textbf{Categor\'ia} & \textbf{Tratamiento} \\
\midrule
1 a 4   & Bajo     & Se acepta y se revisa peri\'odicamente. \\
5 a 9   & Medio    & Requiere mitigaci\'on definida y seguimiento. \\
10 a 14 & Alto     & Requiere mitigaci\'on activa y plan de contingencia probado. \\
15 a 25 & Cr\'itico & Requiere acci\'on inmediata; condiciona decisiones de alcance. \\
\bottomrule
\end{tabularx}

% =====================================================================
\section{Riesgos identificados}

Los riesgos se agrupan seg\'un su origen: los que provienen de las fuentes de
datos y su tratamiento, los que provienen del marco legal aplicable, los que
provienen de dependencias externas y los que afectan la adopci\'on del producto.
Cada ficha registra los ocho campos definidos en la secci\'on anterior.

\subsection*{Riesgos de fuentes y datos}

\begin{riesgo}{Cambios en las fuentes oficiales}
\rdescripcion{Un organismo modifica el formato, la estructura, la ruta o el
mecanismo de acceso de una publicaci\'on sin aviso previo, y el proceso de
ingesta deja de interpretarla correctamente.}
\rcausa{Las fuentes oficiales no asumen compromisos de estabilidad con terceros
que las consumen de forma automatizada. Los rediseños de portales y los cambios
de plataforma de publicaci\'on ocurren sin preaviso.}
\rimpacto{La serie afectada deja de actualizarse o, peor, se actualiza con
valores mal interpretados. En el segundo caso el error se propaga a las
m\'etricas y a los cruces que dependen de ella.}
\rprobabilidad{4 --- Alta}
\rnivel{4 $\times$ 3 = 12 --- Alto}
\rmitigacion{Validaciones autom\'aticas sobre estructura, tipos, rangos y
continuidad de la serie antes de publicar cualquier dato. Ante una desviaci\'on,
la ejecuci\'on se detiene y el registro va a cuarentena en lugar de publicarse.}
\rcontingencia{Se congela la \'ultima versi\'on v\'alida de la serie, se marca
como desactualizada en la interfaz con su fecha de corte, y se adapta el
extractor de esa fuente con prioridad sobre el resto del trabajo.}
\rresponsable{\'Area de procesamiento de datos}
\end{riesgo}

\begin{riesgo}{Discontinuaci\'on o reducci\'on de cobertura de una fuente}
\rdescripcion{Un organismo deja de publicar una serie cubierta por el producto,
reduce su frecuencia o recorta su cobertura hist\'orica.}
\rcausa{Cambios de metodolog\'ia, reasignaci\'on de competencias entre
organismos o decisiones de pol\'itica de publicaci\'on ajenas al producto.}
\rimpacto{Una variable deja de estar disponible para uno o m\'as pa\'ises.
Cuando la variable participa de un cruce, el cruce tampoco puede calcularse.}
\rprobabilidad{2 --- Baja}
\rnivel{2 $\times$ 4 = 8 --- Medio}
\rmitigacion{Conservaci\'on del hist\'orico completo en la capa Bronze, de modo
que lo ya publicado siga disponible aunque la fuente lo retire. Identificaci\'on
de fuentes alternativas para las variables cr\'iticas de cada pa\'is.}
\rcontingencia{Se informa expl\'icitamente la discontinuidad en la ficha de la
variable, se mantiene el hist\'orico accesible con su fecha de corte y se
eval\'ua la fuente alternativa identificada.}
\rresponsable{\'Area de procesamiento de datos}
\end{riesgo}

\begin{riesgo}{Errores o inconsistencias en los datos}
\rdescripcion{Diferencias de fechas, unidades, nombres, horizontes o
metodolog\'ias entre fuentes producen series que parecen comparables sin serlo.}
\rcausa{Cada organismo define sus propias convenciones. Dos series con el mismo
nombre pueden medir cosas distintas, y una misma serie puede cambiar de
metodolog\'ia a lo largo del tiempo.}
\rimpacto{El usuario compara magnitudes no equivalentes y obtiene conclusiones
err\'oneas. Como la comparabilidad es la propuesta central del producto, un
error de este tipo afecta directamente su credibilidad.}
\rprobabilidad{3 --- Media}
\rnivel{3 $\times$ 4 = 12 --- Alto}
\rmitigacion{Normalizaci\'on hacia un modelo com\'un de pa\'is, variable,
unidad, frecuencia y horizonte; control de calidad con env\'io a cuarentena de
los registros dudosos; y registro de la nota metodol\'ogica junto a cada dato,
visible para el usuario.}
\rcontingencia{Se retira la m\'etrica o el cruce afectado, se notifica a los
usuarios que lo consultaron y se publica la correcci\'on con su
justificaci\'on metodol\'ogica en el historial de cambios.}
\rresponsable{\'Area de procesamiento de datos}
\end{riesgo}

\begin{riesgo}{Extracci\'on poco confiable en fuentes no estructuradas}
\rdescripcion{El componente basado en LLM que interpreta HTML y archivos
publicados devuelve valores incorrectos o incompletos sin se\~nalar error.}
\rcausa{Las fuentes no estructuradas no tienen esquema fijo. Un modelo de
lenguaje puede producir una salida bien formada pero equivocada, que no se
distingue de una correcta por su forma.}
\rimpacto{Se incorporan datos err\'oneos con apariencia de v\'alidos, que
alimentan m\'etricas y cruces.}
\rprobabilidad{3 --- Media}
\rnivel{3 $\times$ 3 = 9 --- Medio}
\rmitigacion{La salida del extractor no se publica directamente: atraviesa las
mismas validaciones determin\'isticas que el resto de la ingesta ---tipos,
rangos, continuidad y contraste con la publicaci\'on anterior--- y el archivo
original queda conservado en la capa Bronze para reprocesarlo.}
\rcontingencia{Se reprocesa la fuente afectada desde el original conservado y,
si el problema persiste, se sustituye la extracci\'on autom\'atica por carga
supervisada hasta resolverlo.}
\rresponsable{\'Area de procesamiento de datos}
\end{riesgo}

\subsection*{Riesgos legales}

\begin{riesgo}{Condiciones de uso que no habilitan el uso previsto}
\rdescripcion{Las condiciones de un organismo no permiten almacenar, transformar,
conservar hist\'oricos, exportar o redistribuir sus datos en la forma en que el
producto lo hace.}
\rcausa{No existe una licencia \'unica por pa\'is. Cada organismo fija sus
propias condiciones, que pueden cambiar y que no siempre est\'an expresadas de
forma expl\'icita para el uso automatizado.}
\rimpacto{Obliga a retirar informaci\'on ya publicada, a bloquear un pa\'is o
una fuente completa, y expone al producto a un reclamo del organismo emisor.}
\rprobabilidad{3 --- Media}
\rnivel{3 $\times$ 5 = 15 --- Cr\'itico}
\rmitigacion{Ning\'un conjunto de datos se incorpora sin su entrada en la matriz
de fuentes, que registra instituci\'on, URL, fecha de consulta, condiciones de
uso, atribuci\'on requerida y usos permitidos. La verificaci\'on es fuente por
fuente y se revisa peri\'odicamente. Toda m\'etrica derivada se identifica como
elaboraci\'on propia y conserva la atribuci\'on a la fuente.}
\rcontingencia{Se bloquea de inmediato el conjunto de datos afectado mediante la
funci\'on administrativa prevista, se retira de la interfaz y de las
exportaciones, y se contacta al organismo para solicitar autorizaci\'on expresa.}
\rresponsable{\'Area de an\'alisis de requerimientos}
\end{riesgo}

\subsection*{Riesgos de dependencias externas}

\begin{riesgo}{Dependencia del proveedor externo de cobros}
\rdescripcion{El proveedor de pagos interrumpe su servicio, cambia sus
condiciones o modifica su interfaz de integraci\'on.}
\rcausa{El cobro se delega \'integramente a un tercero para no almacenar ni
procesar datos de tarjetas, lo que traslada esa disponibilidad fuera del
control del producto.}
\rimpacto{No pueden contratarse ni renovarse suscripciones mientras dure la
interrupci\'on. El acceso de los usuarios ya activos no se ve afectado.}
\rprobabilidad{2 --- Baja}
\rnivel{2 $\times$ 3 = 6 --- Medio}
\rmitigacion{Integraci\'on acotada a \emph{checkout} alojado y notificaciones
por \emph{webhook}, sin l\'ogica propia del proveedor dentro del producto, de
modo que sustituirlo afecte una superficie reducida. Los eventos recibidos se
registran para poder reconciliarlos.}
\rcontingencia{Se suspende temporalmente la contrataci\'on en l\'inea
manteniendo activo el acceso vigente, y se reconcilian los estados de
suscripci\'on cuando el servicio se restablece.}
\rresponsable{\'Area de desarrollo backend}
\end{riesgo}

\begin{riesgo}{Indisponibilidad o sobrecosto de la infraestructura}
\rdescripcion{La infraestructura de nube deja de estar disponible dentro de las
condiciones previstas, o su costo supera lo estimado.}
\rcausa{Crecimiento del volumen de datos hist\'oricos, aumento de la frecuencia
de ingesta, o cambios en los precios y las condiciones del proveedor.}
\rimpacto{Degrada la disponibilidad del servicio o presiona la viabilidad
econ\'omica de la operaci\'on.}
\rprobabilidad{2 --- Baja}
\rnivel{2 $\times$ 3 = 6 --- Medio}
\rmitigacion{Despliegue en contenedores desacoplados, que permite dimensionar
cada componente por separado y reduce el costo de migrar. Seguimiento del
consumo real contra lo estimado.}
\rcontingencia{Se reduce la frecuencia de ingesta de las fuentes menos
cr\'iticas y se traslada el hist\'orico frio a almacenamiento de menor costo.}
\rresponsable{\'Area de infraestructura}
\end{riesgo}

\begin{riesgo}{Costo del procesamiento con modelos de lenguaje superior al estimado}
\rdescripcion{El consumo real de los componentes basados en LLM excede la
estimaci\'on utilizada para dimensionar la operaci\'on.}
\rcausa{El costo depende del n\'umero de fuentes, su frecuencia, el tama\~no de
los documentos procesados, los reintentos ante fallas y las validaciones. Ese
volumen todav\'ia no est\'a medido sobre el inventario real de fuentes.}
\rimpacto{Eleva el costo operativo por encima de lo previsto y obliga a revisar
la ecuaci\'on econ\'omica de la operaci\'on.}
\rprobabilidad{3 --- Media}
\rnivel{3 $\times$ 2 = 6 --- Medio}
\rmitigacion{Restringir el uso del modelo a las fuentes que no pueden procesarse
de forma determin\'istica, procesar solo el fragmento relevante de cada
documento, y cachear los resultados por versi\'on de archivo para no reprocesar
lo que no cambi\'o.}
\rcontingencia{Se sustituye el extractor por reglas determin\'isticas en las
fuentes de estructura estable y se reserva el modelo para las restantes.}
\rresponsable{\'Area de procesamiento de datos con LLM}
\end{riesgo}

\subsection*{Riesgos de adopci\'on}

\begin{riesgo}{Baja adopci\'on del p\'ublico objetivo}
\rdescripcion{El producto es percibido como \'util, pero no como
suficientemente necesario para incorporarlo al trabajo habitual.}
\rcausa{El trabajo manual que el producto reemplaza puede estar ya resuelto de
forma aceptable mediante planillas y procesos internos, cuyo costo el usuario
naturaliz\'o.}
\rimpacto{El uso recurrente no se sostiene, y con \'el tampoco la
justificaci\'on de las capacidades construidas sobre el uso peri\'odico.}
\rprobabilidad{3 --- Media}
\rnivel{3 $\times$ 4 = 12 --- Alto}
\rmitigacion{Validaci\'on continua con usuarios reales del segmento objetivo,
medici\'on del uso recurrente sobre las capacidades entregadas y ajuste de las
m\'etricas y funcionalidades seg\'un lo observado.}
\rcontingencia{Se concentra el esfuerzo en el segmento que s\'i muestra uso
recurrente y se revisa la priorizaci\'on de capacidades a partir de esa
evidencia.}
\rresponsable{\'Area de an\'alisis de requerimientos}
\end{riesgo}

% ---------------------------------------------------------------------
%  PLANTILLA: copia este bloque para cada riesgo nuevo.
% ---------------------------------------------------------------------
%  \begin{riesgo}{Titulo del riesgo}
%  \rdescripcion{}
%  \rcausa{}
%  \rimpacto{}
%  \rprobabilidad{}
%  \rnivel{}
%  \rmitigacion{}
%  \rcontingencia{}
%  \rresponsable{}
%  \end{riesgo}

% =====================================================================
\section{Tabla de evaluaci\'on}

La tabla siguiente resume la evaluaci\'on de todos los riesgos registrados,
ordenados por su c\'odigo. El nivel resulta de multiplicar probabilidad por
impacto seg\'un las escalas definidas.

\begin{tabularx}{\textwidth}{@{}l X c c c@{}}
\toprule
\textbf{C\'odigo} & \textbf{Riesgo} & \textbf{Prob.} & \textbf{Impacto} & \textbf{Nivel} \\
\midrule
R-01 & Cambios en las fuentes oficiales                       & 4 & 3 & 12 --- Alto \\
R-02 & Discontinuaci\'on o reducci\'on de cobertura de una fuente & 2 & 4 & 8 --- Medio \\
R-03 & Errores o inconsistencias en los datos                 & 3 & 4 & 12 --- Alto \\
R-04 & Extracci\'on poco confiable en fuentes no estructuradas & 3 & 3 & 9 --- Medio \\
R-05 & Condiciones de uso que no habilitan el uso previsto    & 3 & 5 & 15 --- Cr\'itico \\
R-06 & Dependencia del proveedor externo de cobros           & 2 & 3 & 6 --- Medio \\
R-07 & Indisponibilidad o sobrecosto de la infraestructura    & 2 & 3 & 6 --- Medio \\
R-08 & Costo del procesamiento con modelos de lenguaje        & 3 & 2 & 6 --- Medio \\
R-09 & Baja adopci\'on del p\'ublico objetivo                  & 3 & 4 & 12 --- Alto \\
\bottomrule
\end{tabularx}

% =====================================================================
\section{Matriz de riesgos}

La matriz ubica cada riesgo seg\'un su probabilidad (filas) y su impacto
(columnas), de modo que la concentraci\'on de casos quede visible de un vistazo.
El extremo superior derecho corresponde a los riesgos de mayor exposici\'on.

\begin{tabularx}{\textwidth}{@{}l *{5}{>{\centering\arraybackslash}X}@{}}
\toprule
\textbf{Prob.\textbackslash Impacto} & \textbf{1} & \textbf{2} & \textbf{3} & \textbf{4} & \textbf{5} \\
\midrule
\textbf{5 --- Muy alta} &  &  &  &  &  \\
\addlinespace
\textbf{4 --- Alta}     &  &  & R-01 &  &  \\
\addlinespace
\textbf{3 --- Media}    &  & R-08 & R-04 & R-03, R-09 & R-05 \\
\addlinespace
\textbf{2 --- Baja}     &  &  & R-06, R-07 & R-02 &  \\
\addlinespace
\textbf{1 --- Muy baja} &  &  &  &  &  \\
\bottomrule
\end{tabularx}

\bigskip

De los nueve riesgos registrados, uno es cr\'itico (R-05), tres son altos (R-01,
R-03 y R-09) y cinco son medios. La concentraci\'on en la franja de impacto alto
es coherente con la naturaleza del producto: lo que est\'a en juego en casi
todos los casos es la confiabilidad de la informaci\'on entregada, no la
disponibilidad del servicio.

% =====================================================================
\section{Seguimiento}

El registro se revisa de forma peri\'odica. En cada revisi\'on se verifica si
alg\'un riesgo cambi\'o de probabilidad o impacto, si apareci\'o alguno nuevo y
si las mitigaciones previstas se implementaron efectivamente.

Cada riesgo se encuentra en uno de cuatro estados:

\begin{itemize}
  \item \textbf{Abierto:} identificado, con mitigaci\'on definida pero a\'un no
        implementada por completo.
  \item \textbf{Mitigado:} la mitigaci\'on est\'a implementada y el nivel de
        exposici\'on se redujo respecto de la evaluaci\'on inicial.
  \item \textbf{Materializado:} el evento ocurri\'o; se aplic\'o la
        contingencia y se registra qu\'e efecto tuvo.
  \item \textbf{Cerrado:} la causa desapareci\'o y el riesgo ya no aplica. Se
        conserva en el registro con su fecha de cierre.
\end{itemize}

Las evaluaciones de probabilidad e impacto de la tabla anterior corresponden a
la valoraci\'on vigente y se ajustan cuando la evidencia lo justifica: una
fuente que cambia dos veces de formato eleva la probabilidad de R-01, y una
matriz de fuentes completa y verificada reduce la de R-05.

\end{document}
